\begin{table}[H]
\centering
\caption{Census of Governments municipal debt balance: border sample}
\begin{tabular*}{\textwidth}{@{\extracolsep{\fill}}llccc}
\toprule
Year & Variable & Control & Treat & Difference \\
\midrule
2012 & Long-term debt outstanding & 55.55 & 46.39 & -9.16 \\
 & Total outstanding debt & 58.34 & 47.38 & -10.96 \\
 & Total outstanding debt per capita & 2,467.12 & 1,227.53 & -1,239.58*** \\
 & Population & 23.64 & 22.92 & -0.72 \\
\addlinespace
2017 & Long-term debt outstanding & 57.46 & 50.11 & -7.35 \\
 & Total outstanding debt & 60.59 & 50.65 & -9.94 \\
 & Total outstanding debt per capita & 2,456.77 & 1,451.88 & -1,004.89*** \\
 & Population & 24.07 & 23.56 & -0.51 \\
\addlinespace
2022 & Long-term debt outstanding & 76.09 & 48.11 & -27.97 \\
 & Total outstanding debt & 80.24 & 48.48 & -31.76* \\
 & Total outstanding debt per capita & 3,249.69 & 1,547.13 & -1,702.57*** \\
 & Population & 24.58 & 23.81 & -0.78 \\
\midrule
2012 municipalities &  & 420 & 124 &  \\
2017 municipalities &  & 420 & 124 &  \\
2022 municipalities &  & 420 & 124 &  \\
\bottomrule
\end{tabular*}
\begin{minipage}{\textwidth}
\footnotesize Notes: This table reports means for Census of Governments municipal units matched to the Mergent border-city sample. Treat equals one when the Mergent issuer is in a state requiring city GO bond referendums. Difference is treated minus control. Stars denote Welch two-sample t-tests: * $p<0.10$, ** $p<0.05$, *** $p<0.01$. Long-term debt outstanding is Census item 49U. Total outstanding debt is end-of-year long-term plus short-term debt, in millions of dollars. Population is in thousands.
\end{minipage}
\end{table}
